\section{Reichenbach: bounded calculations and countercases}
\label{sec:annex-p15-reichenbach-probability-logic}

The reference-first notation and boundaries of
Chapter~\ref{sec:p15-reichenbach-probability-logic} apply. These
proofs are not a completeness theorem for the historical calculus
or a proof of its philosophical vindication.

\begin{annexitem}{annex:p15.reichenbach-counting-calculus}{thm:p15-reichenbach-counting-calculus}{Finite counts and qualified limits}
\begin{proof}
On a finite index set put \(P(A,B)=|A\cap B|/|A|\) for
nonempty \(A\). Inclusion of \(A\) in \(A\cup B\),
emptiness of \(B\cap\bar B\), and nonnegativity of counts
prove normalization. Inclusion--exclusion of \(A\cap B\) and
\(A\cap C\) proves addition. For nonempty \(A\cap B\),
\[
 |A\cap B\cap C|/|A|
 =(|A\cap B|/|A|)(|A\cap B\cap C|/|A\cap B|),
\]
proving multiplication. If the branch is empty its joint count is zero;
no conditional ratio on it is evaluated. Partitioning \(A\) by
\(B,\bar B\) gives the complement identity. Partitioning
\(A\cap C\) the same way and applying multiplication on each
nonempty branch gives elimination. The Bayes expression simplifies to
\[
 \frac{(|AB|/|A|)(|ABC|/|AB|)}{|AC|/|A|}
 =|ABC|/|AC|=P(A.C,B),
\]
with all denominators positive by hypothesis.
These equalities hold on every finite initial segment where they are
defined. Sums and products preserve convergence. Division does too
when the denominator has a positive limit: eventually it exceeds half
that limit, and the ratio error is bounded by a constant times the
two errors of numerator and denominator. Thus all the qualified
limiting identities follow, without defining a ratio of zero limits.
\end{proof}
\end{annexitem}

\begin{annexitem}{annex:p15.reichenbach-coupling}{thm:p15-reichenbach-coupling}{Three parameters and every Boolean connection}
\begin{proof}
The four disjoint cells \(ab,a\bar b,\bar a b,\bar a\bar b\)
have masses \(w,u-w,v-w,1-u-v+w\). Their nonnegativity is
equivalent to the bounds stated. Conversely those four nonnegative
numbers sum to one and define a four-point probability model with
the required parameters, proving realizability. Negation collects the
last two cells; disjunction the first three; material implication omits
only \(a\bar b\); equivalence collects the first and fourth.
This gives every entry of the table. For \(u>0\), selection has
degree \(w/u\), hence \(w=ur\), which gives the source's
coupling form and shows why it differs from \(1-u+ur\).
At zero/one marginal values the bounds force \(w=uv\);
substitution gives the four classical assignments. For one die the
faces exclude each other, so \(w=0\) and union \(1/3\).
For two independent dice \(w=1/36\) and union \(11/36\).
Substituting \(2/5,3/5,3/10\) gives \(3/4,7/10,9/10,3/5\)
for coupling, union, material implication and equivalence.
\end{proof}
\end{annexitem}

\begin{annexitem}{annex:p15.reichenbach-frequency-boundaries}{thm:p15-reichenbach-frequency-boundaries}{Closure, sparse selection and frequency one}
\begin{proof}
Take blocks of \(L_k=2^{2^k}\) pairs. In odd blocks use pairs
of coordinate values \((1,1),(0,0)\), and in even blocks
\((1,0),(0,1)\). Each completed pair gives each coordinate
exactly one positive value. The discrepancy at an odd-length prefix
is at most one half, so both marginal frequencies tend to one half.
The conjunction has one positive per odd-block pair and none per
even-block pair. Earlier lengths are negligible because
\(\sum_{j<k}L_j/L_k\leq k2^{-2^{k-1}}\to0\).
Consequently its frequencies at odd-block ends tend to one half and
at even-block ends to zero; no limit exists. At these ends union
frequency is one minus conjunction frequency, so it has no limit
either. This construction does not assume joint normality.

Let \(A\) be the square indices \(j^2\), and let \(B\subset A\)
contain the squares with odd \(j\). Counts through \(n\) are
\(m=\lfloor\sqrt n\rfloor\) and \(\lceil m/2\rceil\).
Both densities and their joint density are zero, while the frequency
within the selected sequence is \(\lceil m/2\rceil/m\to1/2\).
It cannot be obtained as a ratio of the two zero densities. Finally
a series with only its first term false has frequency
\((n-1)/n\to1\), but its all-true universal statement is false.
\end{proof}
\end{annexitem}

\begin{annexitem}{annex:p15.reichenbach-order-and-appraisal}{thm:p15-reichenbach-order-and-appraisal}{Refinement and arbitrary enumeration}
\begin{proof}
In subclass \(j\), selecting the more frequent outcome succeeds
with frequency \(\max(p_j,1-p_j)\). Weighted addition gives
the stated rate. Each maximum exceeds both arguments, so the result
dominates either constant prediction. Equal subclass weights and
rates \(1/4,3/4\) give \(3/4\), whereas the unrefined rate
is \(1/2\).

Enumerate the two infinite kinds separately. For \(0<p<1\),
let \(b_n=\lfloor pn\rfloor-\lfloor p(n-1)\rfloor\).
It is binary, has positive count \(\lfloor pn\rfloor\), and
has infinitely many of each value. At each one take the next unused
positive-class member, otherwise the next unused negative-class member.
Every member is eventually used and the resulting frequency is
\(\lfloor pn\rfloor/n\to p\). For \(p=0\), select
positives at square positions and negatives elsewhere; for \(p=1\)
reverse the kinds. Both lists are exhausted and the limits follow
from the square-count calculation. For no limit, alternate blocks
of next unused positives and negatives, each block longer than the
square of the total earlier length plus one. Both lists are exhausted;
frequencies at successive block ends tend alternately to one and zero.
Thus this is order dependence, not a uniform random draw from an
infinite unordered set.
\end{proof}
\end{annexitem}

\begin{annexitem}{annex:p15.reichenbach-straight-rule-boundary}{thm:p15-reichenbach-straight-rule-boundary}{Attainability, prefix ambiguity and rival rules}
\begin{proof}
The first assertion is the definition of convergence: each positive
tolerance has a tail on which every error is below it. That tail
can depend on the sequence. Given a prefix of length \(m\) with
\(k\) positives, append
\(b_t=\lfloor pt\rfloor-\lfloor p(t-1)\rfloor\), including
constant words at the endpoints. The combined frequency is
\((k+\lfloor pt\rfloor)/(m+t)\to p\), proving that every
prefix admits every limit. For any proposed common \(N\), a
prefix of \(N\) zeros followed by ones has limit one and error
one at \(N\); thus it defeats even tolerance one half.
Clipping a real value to \([0,1]\) cannot increase its distance
from \(p\in[0,1]\): check values below zero, within the interval,
and above one separately. Hence
\[
 |\operatorname{clip}(f_n+d_n)-p|
 \leq|f_n-p|+|d_n|\to0.
\]
Changing finitely many initial predictions also leaves the limit
unchanged. Neither uniqueness nor a finite-stage error estimate follows
from these facts. They are the bounded mathematics behind the original
attainability argument, not a deduction of its epistemic vindication.
\end{proof}
\end{annexitem}

\begin{annexitem}{annex:p15.reichenbach-modal-division}{thm:p15-reichenbach-modal-division}{Modal table and the threshold failure}
\begin{proof}
Let \(S,T\) be the truth sets of the two series in a nonempty
index domain \(I\). They are respectively \(I\), empty, or
proper nonempty sets according to categories \(N,\mathcal J,Pos\).
Negation complements the sets. Product, union, material implication
and equivalence have truth sets
\(S\cap T,S\cup T,(I\setminus S)\cup T\), and
\((S\cap T)\cup(I\setminus(S\cup T))\).
Substituting \(I\) or the empty set proves every row except
Pos/Pos. In that row intersection can be empty or proper nonempty
but never all of \(I\); union and implication can be proper
nonempty or universal but never empty; equivalence can have any of
the three categories. All options are realized: take \(S=T=\{1\}\)
on \(I=\{1,2\}\); take complementary singleton sets on that
domain; take different singletons on \(\{1,2,3\}\); or take
\(\{1,2\},\{2,3\}\) there. This also proves the table's
completeness, not just necessity of its options. In a finite domain,
the three set categories have the indicated counting frequencies.
The preceding frequency-one countercase shows why the reverse
characterization need not hold on an infinite domain.
For threshold \(0<\tau<1\), choose
\(\tau<p<\sqrt\tau\), possible since \(\tau<\sqrt\tau\).
Two independent events with probabilities \(p,p\) exceed the
threshold, while their intersection has probability \(p^2<\tau\).
The numerical example is its direct substitution. Threshold division
therefore is not a truth-preserving Boolean homomorphism.
\end{proof}
\end{annexitem}

\begin{annexitem}{annex:p15.reichenbach-higher-order-boundary}{thm:p15-reichenbach-higher-order-boundary}{A hierarchy is not an unconditional infinite product}
\begin{proof}
Partition \(E\) into \(E\cap H,E\cap\bar H\).
Finite additivity and multiplication give
\(P(E)=qr+(1-q)s\), with both conditionals defined when
\(0<q<1\); at endpoints only the nonnull branch contributes.
The product \(qr\) is the first joint event only. Substitution
gives \(3/8+9/20=33/40\).
For \(0\leq a_j\leq1\), induction on \(m\) gives
\(\prod_{j=1}^m(1-a_j)\geq1-\sum_{j=1}^m a_j\): the step
uses \((1-x)(1-y)=1-x-y+xy\geq1-x-y\), or multiplies
the inductive inequality by the nonnegative next factor. With
\(a_j=2^{-j-2}\), the sum is at most \(1/4\).
The partial products decrease and are bounded below by \(3/4\),
so they have a positive limit. Thus factors strictly below one alone
do not entail collapse. This calculation is not an independence or
coherence proof for the historical hierarchy.
\end{proof}
\end{annexitem}

\begin{annexitem}{annex:p15.reichenbach-finite-fit}{thm:p15-reichenbach-finite-fit}{Distinct global laws with the same finite exact fit}
\begin{proof}
At each \(x_i\), the finite product contains a zero factor, so
every \(g_c(x_i)=L(x_i)\). At any other input every factor is
nonzero, hence
\(g_c(x)-g_d(x)=(c-d)\prod_j(x-x_j)\ne0\) for \(c\ne d\).
All these polynomial laws therefore fit the data but disagree outside
it. Additional finite inputs merely extend the finite product; the
same argument persists. No noisy-data or probability assumption is
needed for this noise-free underdetermination result used in the
historical authors' concluding reasoning case.
\end{proof}
\end{annexitem}
